\section{本章小结}
\label{sec:dual-summary}

本章开头的问题是，手里有一个可行点，怎样知道它离最优还差多少。回答是对偶函数：固定一组合法乘子，对原变量在整个定义域上取下确界，得到的数就是最优值的下界。这个下界只用了下确界的定义和符号，所以弱对偶不需要凸性、可微性和约束资格。对偶函数总是凹的，但计算它的每个值都要做一次内层全局最小化。

一对候选给出的误差证书$f(\hat{\bm{x}})-q(\bm{\lambda},\bm{\nu})$有三个前提：原始可行，乘子合法，$q$算得可信。内层近似得到的值在真正的下界上方，必须扣除经过证明的误差，才能写进证书。

凸共轭把“函数加线性项取下确界”单独提取出来，二次函数、范数、约束指标函数的共轭分别对应二次项、对偶范数球和支撑函数，由此可以直接写出线性约束问题和复合问题的对偶形式。共轭能够重建适当、下半连续的凸函数，但这一点与某个具体问题的零间隙是两回事。

最好的下界是否等于最优值，取决于正则条件。Slater条件和复合问题的相对内部条件都排除了竖直的支撑超平面，保证零间隙和对偶最优乘子存在，却不保证原始最优点存在。凸性本身不能保证零间隙，零间隙也不保证两侧都达到。

实际使用时，能够核对的证据是一对候选加上间隙。零间隙的候选对同时给出两侧的全局最优性，等价于Lagrange鞍点，这些结论不要求凸性；KKT点在凸问题中才等价于它们。乘子如何解释参数的变化，由第~\ref{chap:parametric-sensitivity}章讨论；怎样在算法中实际更新乘子，由第~\ref{chap:penalty-multiplier-methods}章讨论。
